% chapters/02_background.tex
% Drafted 2026-06-07. Readability pass 2026-07-21: connective prose and formal
% definitions re-authored for sequential explanation. Public labels and downstream
% notation are preserved; the mutual-information chain rule and information-
% bottleneck optimization direction are corrected.
% Shared formal vocabulary, per outlines/02_background.md.
% Two boundaries: (i) primitives-only — named phenomena/methods are defined where
% contributed (Label-Blindness=Ch.4, DSC/TGT=Ch.5, layer-pair probe=Ch.6), so v1's
% "Domain Feature Collapse" definition is deleted here and replaced by a DSC
% forward-pointer; (ii) Ch.2 DEFINES the primitive, Ch.3 surveys its use. Sections
% run tasks (2.1-2.5, ported from v1 §2) then the two lenses (2.6-2.8) in the order
% Ch.3 opens with them. Notation follows math_commands.tex; geometry matches Ch.5,
% the InfoNCE/DPI primitives match Appendix C. Citations are sparse by design.
\chapter{Background and Definitions}
\label{ch:background}

This chapter defines the concepts used by more than one results chapter. The first
five sections specify the tasks and modeling settings. The final three introduce
information theory and representation geometry. Named results and methods remain
in the chapters that contribute them
(\Cref{ch:label-blindness,ch:domain-sensitivity-collapse,ch:hallucination-detection}).
\Cref{ch:literature-review} reviews how prior work uses the same concepts. Citations
here are therefore limited to foundational attributions.

\section{Out-of-Distribution Detection}
\label{sec:bg:ood}

Out-of-distribution (OOD) detection asks whether an input belongs outside the
model's training distribution. This dissertation focuses on \emph{semantic
shift}~\citep{yang2021generalized}. A semantic shift occurs when the true label is
not among the in-distribution labels, even if the input resembles the training
data.

\begin{definition}[Out-of-Distribution Detection]
Given an input $\rvx$ and a set of in-distribution labels $\sY_{\text{in}}$, OOD
detection determines whether the true label $\rvy$ lies outside that set:
\[
  \rvy \notin \sY_{\text{in}}.
\]
Equivalently, a probabilistic detector estimates
$P(\rvy \notin \sY_{\text{in}} \mid \rvx)$.
\label{def:bg:ood}
\end{definition}

The canonical baseline reads a confidence signal directly off a trained
classifier's output distribution.

\begin{method}[Maximum Softmax Probability (MSP)]
Given an input $\rvx$, the MSP method defines its OOD score as
$1-\mathrm{MSP}(\rvx)$, where
$\mathrm{MSP}(\rvx) = \max_{y \in \sY_{\text{in}}} P(y \mid \rvx)$~\citep{hendrycks2016baseline}.
\label{method:bg:msp}
\end{method}

We assume that a labeling function $f_{\rvy}$ assigns labels from the input alone.
Let $\sY_{\text{all}}$ contain every label that this function can produce over the
full input space. A training set need not contain them all, so in general
$\sY_{\text{in}} \subsetneq \sY_{\text{all}}$.

OOD methods differ in how they interact with training. A
\emph{training-agnostic} method, such as MSP, scores an existing classifier without
changing it. A \emph{training-aware} method changes the architecture, objective, or
augmentation procedure to improve detection.

OOD detection primarily addresses \emph{epistemic} uncertainty caused by limited
knowledge. It is distinct from \emph{aleatoric} uncertainty caused by irreducible
noise in the data. \Cref{ch:label-blindness,ch:domain-sensitivity-collapse} use
this OOD setting.

\section{Anomaly Detection}
\label{sec:bg:anomaly}

Anomaly detection and OOD detection both reject unusual inputs, but they use
different notions of unusual. \Cref{tab:bg:anomaly-ood} summarizes the distinction.

\begin{table}[H]
  \centering
  \caption{Typical distinctions between anomaly detection and OOD detection.}
  \label{tab:bg:anomaly-ood}
  \begin{tabular}{@{}lll@{}}
    \toprule
    & \textbf{Anomaly detection} & \textbf{OOD detection} \\
    \midrule
    Target
      & \shortstack[l]{Atypical points within\\a modeled domain}
      & \shortstack[l]{Inputs outside the modeled\\label set or distribution} \\
    Training data
      & \shortstack[l]{Usually normal examples\\from one class}
      & \shortstack[l]{Usually examples from several\\in-distribution classes} \\
    Evaluation
      & \shortstack[l]{Separates normal points\\from rare deviations}
      & \shortstack[l]{Separates in-distribution from\\out-of-distribution inputs} \\
    \bottomrule
  \end{tabular}
\end{table}

Anomaly detection is therefore often a \emph{one-class} problem: training provides
normal examples but no examples of the anomalies that must later be rejected. The
same training pattern appears in the unlabeled detectors of
\Cref{ch:label-blindness}. The hallucination probe of
\Cref{ch:hallucination-detection} uses a related but supervised asymmetric
construction. Its training data contain both faithful and hallucinated
generations, but only faithful generations form a mutual-positive class;
hallucinated generations otherwise act as negatives while retaining their own
cross-layer view pairs as positives.

\section{Unlabeled Out-of-Distribution Detection}
\label{sec:bg:unlabeled}

Most OOD detectors are trained on labeled in-distribution data and repurpose a
supervised classifier's outputs, softmax confidence or logits, as the detection
signal. Not all do. We call a detector \emph{unlabeled} when its training uses only
the inputs and never the labels.

\begin{definition}[Unlabeled OOD Detection]
Let the in-distribution dataset be
\[
  \sD_{\text{in}}=\{(\rvx_i,\rvy_i)\}_{i=1}^{N}.
\]
An OOD detector is \emph{unlabeled} if training uses only the inputs. It produces
\[
  f_{\text{OOD}}\gets\mathrm{Train}(\{\rvx_i\}_{i=1}^{N})
\]
with no dependence on the labels $\{\rvy_i\}$.
\label{def:bg:unlabeled}
\end{definition}

Unlabeled detectors commonly use density estimation, reconstruction error, or
self-supervised representations. A pretrained model used without fine-tuning also
fits the definition because it never trains on the in-distribution labels.

These methods are useful when labels are expensive or when detection should not
depend on one classification task. They still estimate the same event as a labeled
detector, $P(f_{\rvy}(\rvx) \notin \sY_{\text{in}})$. \Cref{ch:label-blindness}
studies the limit created by removing labels from training.

\section{Dataset Domain and the Single-Domain Setting}
\label{sec:bg:domain}

A \emph{domain} records the conditions under which data are produced, such as the
environment, sensor, or acquisition process. A domain-labeling function
$f_{\rvd}$ assigns each input a domain label,
$\rvd = f_{\rvd}(\rvx)$.

In the single-domain setting, every training input shares one domain $\rvd_1$:
\[
  f_{\rvd}(\rvx) = \rvd_1
  \qquad \text{for every in-distribution } \rvx.
\]
An input for which $f_{\rvd}(\rvx) \neq \rvd_1$ comes from a different domain.

For analysis, write the input as $\rvx=(\rvx_{\rvd},\rvx_{\rvy})$. The domain
features $\rvx_{\rvd}$ identify the generating conditions, while the class features
$\rvx_{\rvy}$ identify the target class. We assume that these feature sets are
independent within the training distribution:
\[
  I(\rvx_{\rvd};\rvx_{\rvy})=0.
\]
This independence is a modeling assumption, not a claim about the full input
space. Outside the training distribution, domain cues may also carry class
information.

\begin{definition}[Single-Domain Dataset]
A dataset is \emph{single-domain} with respect to $f_{\rvd}$ if every training
input has the same domain label $\rvd_1$. A multi-domain dataset contains training
inputs with more than one domain label. For the single-domain analyses in this
dissertation, the input admits domain and class features
$(\rvx_{\rvd},\rvx_{\rvy})$ that satisfy
$I(\rvx_{\rvd};\rvx_{\rvy})=0$ within the training distribution.
\label{def:bg:singledomain}
\end{definition}

Because the training domain is constant, class supervision does not directly
reward sensitivity to domain variation. \Cref{ch:domain-sensitivity-collapse}
studies how this setting can suppress domain-sensitive directions and how training
can restore them. \Cref{sec:bg:geometry} defines the geometric vocabulary used in
that analysis.

\section{Large Language Models and Hallucination}
\label{sec:bg:llm}

The hallucination study uses autoregressive language models and their intermediate
activations. This section defines both objects before introducing the
information-theoretic tools used to analyze them.

\begin{definition}[Large Language Model]
A \emph{large language model} (LLM) is a parameterized distribution over token
sequences. Let $\mathcal{X}$ be a finite vocabulary and $\mathcal{X}^{*}$ the set
of finite sequences over it. For a sequence
$\vx=(x_1,\ldots,x_T)$, an autoregressive model factorizes
\[
  \Pr_{\vtheta}(\vx)
  = \prod_{t=1}^{T}\Pr_{\vtheta}(x_t\mid x_{<t}),
  \qquad x_{<t}=(x_1,\ldots,x_{t-1}).
\]
The distribution is represented by a deep network, typically a Transformer. It is
trained on a corpus $\mathcal{D}$ by minimizing next-token cross-entropy:
\[
  \E_{\vx\sim\mathcal{D}}
  \left[-\sum_t\log\Pr_{\vtheta}(x_t\mid x_{<t})\right].
\]
\label{def:bg:llm}
\end{definition}

The term ``large'' refers to model and dataset scale rather than a fixed numerical
threshold. The analysis below depends on the autoregressive structure and exposed
activations, not on a particular parameter count.

\begin{definition}[Hallucination]
Let $\vx$ be a prompt and $\vy$ a continuation produced by the model. Fix a
reference world model $\mathcal{W}$ that represents the facts and sources
available for evaluation. A \emph{hallucination} occurs when a span
$\vy'\subseteq\vy$ presents as fact content that contradicts $\mathcal{W}$ or
cannot be verified from it.
\label{def:bg:hallucination}
\end{definition}

An \emph{extrinsic} hallucination conflicts with or lacks support from the external
world model $\mathcal{W}$. An \emph{intrinsic} hallucination is inconsistent with
the prompt or with earlier tokens in the generation. Because grounding is relative
to $\mathcal{W}$, an evaluation must state which facts and sources it treats as
authoritative.

\Cref{ch:hallucination-detection} does not assign its benchmark regime to either
category. It uses an operational label: a greedy generation is labeled
hallucinated when the benchmark's gold-reference evaluator marks it incorrect.
This convention fixes the detection target consistently without claiming that
every generation marked incorrect belongs to one hallucination taxonomy.

\paragraph{The residual stream.}
A Transformer maintains a hidden state at every token position and layer. For one
position, write the sequence as $h_1,\ldots,h_L$, where
$h_\ell\in\R^d$ is the state produced by layer $\ell$. This sequence is the
\emph{residual stream}.

The activations are available from one forward pass and require no resampling. An
internal-state detector can therefore score them directly.
\Cref{ch:hallucination-detection} learns from pairs of layers. The next two
sections define the information-theoretic tools that justify this pairing.

\section{Information Theory}
\label{sec:bg:infotheory}

Information theory~\citep{shannon1948mathematical} quantifies uncertainty and
statistical dependence. In this dissertation, it measures how much information a
representation retains about a target.

\begin{definition}[Entropy and Conditional Entropy]
The \emph{entropy} $H(\rvx)$ of a random variable $\rvx$ with distribution
$p(\vx)$ is its expected uncertainty. For discrete variables,
\[
  H(\rvx)=-\sum_{\vx}p(\vx)\log p(\vx).
\]
The \emph{conditional entropy} is the uncertainty remaining in $\rvy$ after
$\rvx$ is known:
\[
  H(\rvy\mid\rvx)
  =-\sum_{\vx,\vy}p(\vx,\vy)\log p(\vy\mid\vx).
\]
The corresponding definitions for continuous variables replace sums with
integrals. Entropy and conditional entropy satisfy the chain rule
\[
  H(\rvx,\rvy)=H(\rvx)+H(\rvy\mid\rvx).
\]
\label{def:bg:entropy}
\end{definition}

\begin{definition}[Mutual Information]
The \emph{mutual information} $I(\rvx;\rvy)$ is the reduction in uncertainty about
one variable after observing the other. Equivalently, it is the divergence between
the joint distribution and the product of the marginals:
\[
  I(\rvx;\rvy)
  =D_{\mathrm{KL}}\!\left(p(\vx,\vy)\,\middle\|\,p(\vx)p(\vy)\right).
\]
Here the Kullback--Leibler divergence is
\[
  D_{\mathrm{KL}}(p\|q)
  =\E_{\vx\sim p}\!\left[\log\frac{p(\vx)}{q(\vx)}\right].
\]
Mutual information is non-negative and equals zero exactly when the variables are
independent. Its chain rule is
\[
  I(\rvx;\rvy,\rvz)
  =I(\rvx;\rvz)+I(\rvx;\rvy\mid\rvz).
\]
\label{def:bg:mi}
\end{definition}

Processing cannot increase the information one variable contains about another.
The data processing inequality states this constraint.

\begin{definition}[Data Processing Inequality]
Suppose three random variables form a Markov chain
$\rvx\rightarrow\rvy\rightarrow\rvz$. Then $\rvz$ depends on $\rvx$ only through
$\rvy$, and
\[
  I(\rvx;\rvz)\le I(\rvx;\rvy).
\]
No transformation of $\rvy$ can increase its information about $\rvx$.
\label{def:bg:dpi}
\end{definition}

The inequality also applies when each of two variables is mapped to a learned
embedding. If $g_a$ and $g_b$ encode two layer activations, then
\[
  I(g_a(h_a);g_b(h_b))\le I(h_a;h_b).
\]
A positive lower bound at the embedding level therefore certifies positive
dependence between the underlying activations. \Cref{ch:hallucination-detection}
uses this fact, and \Cref{app:mi-hallucination-proofs} gives the formal statement.

\section{Sufficiency, Information Bottleneck, and InfoNCE}
\label{sec:bg:ib}

Three concepts turn the information-theoretic view into a tool for representation
learning. Sufficiency states what a representation must retain. The information
bottleneck adds a compression objective. InfoNCE supplies a contrastive lower bound
that can be estimated from samples.

\begin{definition}[Minimal Sufficient Statistic]
A statistic $T(\rvx)$ is \emph{sufficient} for $\rvy$ if it retains all the
information that $\rvx$ carries about $\rvy$:
\[
  I(\rvx;\rvy\mid T(\rvx))=0.
\]
It is \emph{minimal} if it is a function of every other sufficient statistic. A
minimal sufficient statistic is therefore the most compact representation that
loses no information about $\rvy$.
\label{def:bg:mss}
\end{definition}

\begin{definition}[Information Bottleneck]
Given an input $\rvx$ and target $\rvy$, the \emph{information
bottleneck}~\citep{tishby2000information} seeks a representation $\rvz$ that
retains target information while compressing the input. One minimization form is
\[
  \min_{p(\vz\mid\vx)}\;
  \mathcal{L}_{\mathrm{IB}}
  = I(\rvz;\rvx)-\beta I(\rvz;\rvy),
\]
where $\beta>0$ controls the tradeoff between compression and target retention.
Under appropriate capacity and optimization assumptions, this objective favors a
minimal sufficient representation of the target.
\label{def:bg:ib}
\end{definition}

Sufficiency drives the result of \Cref{ch:label-blindness}. If the surrogate task
is independent of the class-relevant features, its minimal sufficient
representation contains no class information. A detector built only on that
representation cannot distinguish the held-out classes.

The language-model analysis asks a different question: whether information remains
between two layer representations. Because this quantity must be estimated from
samples, the analysis uses a contrastive lower bound.

\begin{method}[InfoNCE Lower Bound]
Under the standard InfoNCE sampling construction, let two encoded views
$\mathrm{view}_a$ and $\mathrm{view}_b$ be scored against $K-1$ negatives. If
$\mathcal{L}_{\mathrm{NCE}}$ is the resulting loss~\citep{oord2018representation},
then
\[
  I(\mathrm{view}_a;\mathrm{view}_b)
  \ge \log K-\mathcal{L}_{\mathrm{NCE}}.
\]
Thus $\mathcal{L}_{\mathrm{NCE}}<\log K$ gives a strictly positive lower bound on
the mutual information between the views.
\label{method:bg:infonce}
\end{method}

InfoNCE underlies modern contrastive representation learning. Self-supervised
methods apply it to augmented views~\citep{chen2020simclr}. Supervised contrastive
learning instead treats examples from the same class as
positives~\citep{khosla2020supervised}.

\Cref{ch:hallucination-detection} starts from this construction but uses an
asymmetric supervised contrastive objective: faithful samples form a
mutual-positive class, whereas each hallucinated sample retains only its own
cross-layer views as positives. Because this is a multi-positive supervised
specialization, the standard $\log K$ threshold above does not transfer unchanged.
The class-conditioned bound for the chapter's exact objective appears in
\Cref{app:mi-hallucination-proofs}.

\section{Representation Geometry}
\label{sec:bg:geometry}

Representation geometry asks how learned features are arranged in space. Of
particular interest is how variance is distributed across feature directions.

Let $z=f_\theta(\rvx)\in\R^d$ be the feature assigned to an input from a
$C$-class dataset. Let $\Sigma$ be the empirical covariance of these features,
with eigenvalues $\lambda_1\ge\cdots\ge\lambda_d$. The covariance decomposes as
\[
  \Sigma=\Sigma_{\mathrm{between}}+\Sigma_{\mathrm{within}}.
\]
The between-class component has rank at most $C-1$. If training drives
$\Sigma_{\mathrm{within}}$ toward zero, the representation concentrates its
variance in at most $C-1$ directions.

\begin{definition}[Anisotropy: Participation Ratio and Effective Rank]
The \emph{participation ratio} and \emph{effective rank} measure the concentration
of a covariance spectrum:
\[
  \mathrm{PR}(\Sigma)
  =\frac{(\sum_i\lambda_i)^2}{\sum_i\lambda_i^2},
  \qquad
  r_{\mathrm{eff}}(\Sigma)
  =\exp\!\left(-\sum_i\hat{\lambda}_i\log\hat{\lambda}_i\right),
\]
where $\hat{\lambda}_i=\lambda_i/\sum_j\lambda_j$. Both quantities estimate how
many directions carry meaningful variance. A representation with
$r_{\mathrm{eff}}\ll d$ is strongly anisotropic.
\label{def:bg:anisotropy}
\end{definition}

\begin{definition}[Class-Discriminative Subspace]
Let $v_1,\ldots,v_k$ be the top eigenvectors of $\Sigma$. In the supervised
settings studied here, these directions capture the dominant between-class
variation. We therefore write the \emph{class-discriminative subspace} as
\[
  \mathcal{S}_{\mathrm{cls}}
  =\operatorname{span}\{v_1,\ldots,v_k\}.
\]
Let $P$ be the orthogonal projector onto $\mathcal{S}_{\mathrm{cls}}$ and let
$P_\perp=I-P$. Every feature decomposes as
\[
  z=z_\parallel+z_\perp,
  \qquad z_\parallel=Pz,
  \qquad z_\perp=P_\perp z.
\]
The term ``class-discriminative'' records the supervised alignment assumption. If
that alignment does not hold, the same construction is simply the dominant
principal subspace. A domain shift with little in-distribution variance may appear
primarily in the complementary directions $z_\perp$.
\label{def:bg:clssubspace}
\end{definition}

Detection scores respond differently to this geometry. Euclidean
nearest-neighbor scores~\citep{sun2022out} are often dominated by high-variance
directions. A regularized Mahalanobis score~\citep{lee2018simple} can amplify a
low-variance direction only if the representation preserves useful separation
there.

Logit-based scores read features through the classifier head. Maximum softmax
probability (\Cref{method:bg:msp}) and the energy score~\citep{liu2020energy}
cannot respond to a direction that lies in the head's null space.
\Cref{ch:domain-sensitivity-collapse} makes these conditions precise. Neural
collapse describes the limiting geometry that motivates the analysis.

\begin{definition}[Neural Collapse]
\emph{Neural collapse}~\citep{Papyan2020nc} is a terminal-phase classification
geometry. Within-class variance vanishes,
$\Sigma_{\mathrm{within}}\to0$, and the class means form a
$(C-1)$-simplex equiangular tight frame. The remaining $d-(C-1)$ directions carry
zero variance in the limiting representation.
\label{def:bg:neuralcollapse}
\end{definition}

Neural collapse is an asymptotic limit. When shift-sensitive directions lie in the
collapsed complement, a detector cannot use them through the learned
representation. \Cref{ch:domain-sensitivity-collapse} studies the corresponding
finite-training geometry and a remedy. The formal results appear in
\Cref{app:dsc-proofs}.

The chapter has now defined the two analytical perspectives used in the
dissertation. \Cref{ch:literature-review} reviews their use in prior work, and the
three results chapters apply them to their respective detection problems.
